\chapter{Connection to failure events and maintenance}
\label{ch:failure}

The preceding chapters established that motion defines four operational states, that the
air-system sensors independently recover the deceleration state, and that they explain a
moderate share of braking intensity. This chapter turns to the question behind predictive
maintenance: does braking behaviour change in the days before a documented brake-system failure?
The evidence points to a near-null answer, reported here as a genuine negative rather than
qualified away.

\section{Setup}

The training window contains 10 documented failures (8 brake-system) and 11 maintenance events.
Deceleration events in the seven days before any brake failure form a pre-failure group (about
23{,}000 events), compared against a baseline more than 30 days from any failure (about 44{,}000).
Ten event-level features are tested, spanning the kinematic intensity measures and the auxiliary
health channels.

\section{Distributional test}
\label{sec:prefailure-test}

Each feature is compared between the pre-failure and baseline groups with a Mann--Whitney $U$
test~\cite{MannWhitney:47}, a rank-based test that asks whether values from one group tend to exceed the other without
assuming normality: it ranks the pooled values and counts, over every cross-group pair, how often
the pre-failure value is larger. Under the null that count is about $n_1 n_2 / 2$, and the
$p$-value measures the departure from it. With ten features tested, the $p$-values are
Bonferroni-corrected: at $m=10$ and $\alpha = 0.05$ a feature counts as significant only if its
raw $p$ falls below $\alpha/m = 0.005$, capping the family-wide false-positive chance at $\alpha$.
With more than 23{,}000 pooled events almost every feature clears that bar, which means little
here: at this sample size even trivial differences drive $p$ to zero. What is read instead is the
effect size, Cliff's delta~\cite{Cliff:93}, where $0$ is no separation and $\pm 1$ total, with $|\delta| < 0.15$
treated as negligible.

The result is near-null (Table~\ref{tab:prefailure}). The kinematic and braking-command features
(peak and mean deceleration, deceleration energy, brake-cylinder and pneumatic-braking-force
channels) all sit at $|\delta| < 0.06$, effectively no separation despite significant $p$-values.
Only energy-braking-resistance clears the threshold, at $|\delta| \approx 0.20$; the next largest,
load pressure at $\approx 0.10$, stays negligible. Even that effect is small and rests on only
eight brake failures, so it cannot support a prediction claim on its own. It does hint at
\emph{where} any signal might live, in the auxiliary channels rather than the deceleration
kinematics, a point the conclusion returns to.

\begin{table}[htbp]
  \centering
  \caption{Pre-failure versus baseline effect sizes (Cliff's delta), top 7 $|\delta|$
  values. The braking-kinematic features are negligible; only one auxiliary channel
  (energy-braking-resistance) rises above $|\delta| = 0.15$, and even that is small and rests
  on eight brake failures. Negative $\delta$ means lower values before failure.}
  \label{tab:prefailure}
  \begin{tabular}{lrc}
    \toprule
    Feature & Cliff's $\delta$ & Above negligible ($|\delta|>0.15$)? \\
    \midrule
    energy-braking-resistance (mean)   & $-0.201$ & yes \\
    load pressure (mean)               & $-0.103$ & no  \\
    jerk RMS                           & $+0.054$ & no  \\
    brake-cylinder pressure (integral) & $-0.040$ & no  \\
    deceleration energy (specific)     & $-0.039$ & no  \\
    pneumatic braking force (mean)     & $-0.032$ & no  \\
    peak deceleration                  & $-0.005$ & no  \\
    \bottomrule
  \end{tabular}
\end{table}

\section{Change-point detection}
\label{sec:changepoint}

As an independent check, a CUSUM (cumulative-sum) change-point detector~\cite{Page:54} was run on the strongest
signal from the previous section, the weekly mean of energy-braking-resistance
(Figure~\ref{fig:cusum}). One design choice matters here and applies to this figure alone. The raw
weekly proxy climbs almost monotonically across the year, which would only make sense if brake wear
never reversed. Assuming the brakes are serviced, the proxy is reset to its post-service level at
every maintenance date and only the within-interval drift is read. This is an assumption, not a
fact: the log records that a stop happened, not that the brakes were touched, so the reset may not
hold at every stop. It is adopted because a proxy that only ever rises is not sensible to run a
change-point test on.

CUSUM keeps a running total of each week's deviation from the established mean: while the process
is stable the positive and negative deviations cancel and the sum stays near zero, but once the
level shifts they share a sign and the sum climbs until it crosses a threshold and flags a
change-point. A persistent drift thus surfaces where one-off noise does not. The threshold is four
standard deviations, tuned for about one false alarm every six months. On the reset signal it
fires four times, only one within seven days of a documented failure. There is no reliable
alignment between change-points and failures during the deceleration phase, consistent with the
near-null effect sizes of Section~\ref{sec:prefailure-test}.

\begin{figure}[htbp]
  \centering
  \includegraphics[width=0.95\linewidth]{prefailure_cusum.png}
  \caption{CUSUM change-points on the weekly mean of energy-braking-resistance, the auxiliary
  channel with the largest pre-failure effect size, reset to its post-service level at each
  maintenance date (dotted) under the assumption that the brakes are serviced there. Within
  each between-maintenance interval the proxy drifts up from zero; four change-points fire and
  only one falls within seven days of a documented failure, so even the strongest proxy does
  not track failures during the deceleration phase.}
  \label{fig:cusum}
\end{figure}

\section{Energy and the maintenance calendar}

The deceleration-energy axis of Section~\ref{sec:energy} was examined descriptively against the
maintenance calendar (Figure~\ref{fig:maint}). Energy totals track how long an interval ran and how
much the train moved, following interval length more than failure risk. A CUSUM on weekly summed
energy fires three times, again only one within seven days of a failure. As an accounting check,
about 96\% of all real-deceleration energy falls inside the between-maintenance intervals, so
almost no braking is missed.

One pattern in the same figure is worth flagging, with care. Every between-maintenance interval
containing a brake failure ran longer than about three weeks, and every shorter one was
failure-free: the groups separate cleanly at roughly 21 days. On the service-at-each-stop reading
of Section~\ref{sec:changepoint}, a longer unmaintained stretch gives wear more time to build,
which would make a cadence under three weeks a sensible target. Two cautions hold. A longer
interval also spans more calendar days, so it has more chance to contain any event at all,
confounding length with exposure; and the pattern rests on ten intervals and eight brake failures.
It is an idea worth testing on more data, not a result.

\begin{figure}[htbp]
  \centering
  \includegraphics[width=0.9\linewidth]{maintenance_interval_energy.png}
  \caption{Deceleration energy per between-maintenance interval. Totals rise with interval
  length. Every interval that contained a brake failure (red) ran longer than about three
  weeks; every shorter interval was failure-free. The split is suggestive but confounded with
  exposure time and rests on ten intervals.}
  \label{fig:maint}
\end{figure}
